\section{Cost details and reference limitations}\label{app:cost}
The archived run identifiers in Table~\ref{tab:cost} are \texttt{gpt-oss-120b}, \texttt{gemini-3.5-flash-lite}, and \texttt{gemini-3.5-flash}, including window and verification variants. These names and rates are reported as recorded experiment metadata; they are not a current service catalog. The cost script records input/output prices per million tokens of \$0.04/\$0.17 for OSS, \$0.30/\$2.50 for Lite, and \$1.50/\$9.00 for Flash. The results file documents a historical pricing correction. No present-day price verification is implied by quoting the dated experiment.

For self-verification, the rounded 23,000-page total decomposes into \$22.65 extraction and approximately \$4.17 verification. The implied per-1,000-page components are approximately \$0.985 and \$0.181, respectively. For the Lite and Flash verifier combinations, the corresponding incremental projections are approximately \$0.907 and \$3.599 per 1,000 pages. These calculations use the same recorded extraction stage and should not be summed across alternative pipelines as if each comparison necessarily incurred a fresh extraction bill.

The reference's small size limits precision. One additional matched claim changes recall by $1/52$, or about 1.9 percentage points. Claims within a page share language and layout, and the hard cases were chosen to probe boundaries rather than estimate their population frequency. Conventional claim-independent error bars would obscure these design limitations. The original protocol calls for independent human recoding before publication-level validation; that recoding has not been established by the separate ProQuest audit.

The development record also reports a truncation failure in an early extractor run: a response could exhaust its generation budget and return incomplete structured output. Correcting that behavior changed the measured extraction ranking. This is a useful engineering lesson about auditing failed and truncated calls, but the reported post-fix metrics are not a randomized experiment identifying the effect of token budget. Similarly, comparisons that change both reasoning effort and window size cannot attribute a performance difference to either factor alone.

